% !PRINTED 233
\weberpage{233}

Since the left side is not changed by the permutation of $\sigma$ with $-\sigma$, it is a pure quadratic equation, and, having regard to \S.\,73, (12), it gives
\[
B_{0,3}^2\,\sigma^2 - 5\,D = 0,
\]
where $D$ is the discriminant of the given equation. By (2) we therefore obtain
\weq{12}{5\,a_0\,a_5\,\sigma = \sqrt{5\,D},}
in which the sign is arbitrary.

\begin{center}
\S.\ 75.

\gesperrt{Normal form of the equation of the fifth degree.}
\end{center}

As already remarked earlier, the principal aim of these considerations is to establish a normal form of the equation of the fifth degree which depends on only \gesperrt{one} undetermined coefficient, \gesperrt{one parameter}. Such a normal form is the \gesperrt{Bring-Jerrard} form. To obtain it, by (4) of the preceding paragraph we have to set
\weq{1}{y = t_3\,F_0 + t_2\,(\alpha_2\,F_1 + \alpha_1\,F_2 + \alpha_0\,F_3)}
so that identically
\[
S(y) = 0, \quad S(y^2) = 0
\]
holds, and then the ratio $t_3 : t_2$ is to be formed from the cubic equation
\[
S(y^3) = 0
\]
This cubic equation can actually be formed, even though its expression becomes long. The formulae needed for this have been computed by \gesperrt{Cayley}\textsuperscript{1)}.

It is the result of a remarkable investigation of \gesperrt{Gordan} that one can obtain another normal form, the \gesperrt{Brioschi} form, without any new irrationality, and to conclude these considerations on the Tschirnhausen transformation we shall now derive this result\textsuperscript{2)}.

We hold fast to the assumptions of the preceding paragraph and, for the moment, set as an abbreviation
%CONT>>

\bigskip
\noindent\rule{4cm}{0.4pt}

\noindent\textsuperscript{1)} \gesperrt{Cayley}, on Tschirnhausen's Transformation, Phil. Trans. 1861, Mathematical Papers, Tom. IV, Nr. 275.

\noindent\textsuperscript{2)} \gesperrt{Gordan}, Mathematische Annalen, Bd. 28, 1886.
